% fig_spectral_gap.tex
% Three-panel: (a) spectral gap Delta_r = sigma_r/sigma_{r+1} vs time after fault
%              (b) coupling ratio rho(t) = ||H_a||/||H_d|| vs time
%              (c) bound C_rho = (1+rho)/(1-rho) and resulting kappa bound
\begin{tikzpicture}
\begin{groupplot}[
  group style={group size=1 by 3, vertical sep=0.35cm},
  thesis style,
  width  = 0.90\columnwidth,
  height = 2.6cm,
  xmin=0, xmax=200,
  xtick={0,40,80,120,160,200},
  grid=both,
  grid style={gray!20,line width=0.3pt},
  clip=true,
]

%% ── Panel 1: Spectral gap Delta_r = sigma_r / sigma_{r+1} ───────────────────
\nextgroupplot[
  ylabel={$\Delta_r = \sigma_r/\sigma_{r+1}$},
  ymin=1.0, ymax=8.0,
  ytick={1,2,3,4,5,6,7,8},
  xticklabels={},
  legend pos=south east,
  legend style={font=\scriptsize},
]
%% Spectral gap: rises sharply in subtransient (model well-separated),
%% then decreases as damper windings decay
\addplot[thick,thesisblue,mark=none,domain=0:200,samples=80]
  { (x<5) ? 1.2 :
    (x<40) ? 1.2 + 5.8*(1-exp(-(x-5)/15)) :
    max(2.1, 7.0 - (x-40)*0.028) };
\addlegendentry{$r{=}2$, $n{=}6$ (IBR fault model)};

%% Threshold: gap > 1 always (singular values distinct)
\addplot[thesisred,dotted,thick,domain=0:200,samples=2] {1.0};
\node[thesisred,font=\tiny,anchor=west] at (axis cs:202,1.2) {$\Delta_r{=}1$};

%% Bound-tight region (large gap → bound tight)
\draw[thesisgreen!60,fill=thesisgreen!10]
  (axis cs:10,4.5) rectangle (axis cs:150,8.0);
\node[thesisgreen,font=\tiny] at (axis cs:80,7.5) {Bound tight ($\Delta_r > 4$)};

%% ── Panel 2: Coupling ratio rho ─────────────────────────────────────────────
\nextgroupplot[
  ylabel={$\rho = \|\mathbf{H}_a\mathbf{Q}\|/\|\mathbf{H}_d\|$},
  ymin=0, ymax=0.22,
  ytick={0,0.05,0.10,0.15,0.20},
  xticklabels={},
]
%% Coupling ratio: small initially (algebraic states decouple at subtransient),
%% rises as transient develops
\addplot[thick,thesisorange,mark=none,domain=0:200,samples=80]
  { (x<10) ? 0.02 :
    (x<80) ? 0.02 + (x-10)*0.00085 :
    min(0.19, 0.08 + (x-80)*0.00092) };

%% rho < 0.10 region (contraction condition)
\addplot[thesisgreen,dotted,thick,domain=0:200,samples=2] {0.10};
\node[thesisgreen,font=\tiny,anchor=west] at (axis cs:202,0.105)
  {$\rho_{\max}{=}0.10$};

\draw[thesisgreen!60,fill=thesisgreen!10]
  (axis cs:0,0) rectangle (axis cs:93,0.10);
\node[thesisgreen,font=\tiny] at (axis cs:47,0.06)
  {$\rho<0.10$};

%% ── Panel 3: Resulting kappa bound ──────────────────────────────────────────
\nextgroupplot[
  ylabel={$\kappa$ bound},
  xlabel={Time after fault inception (ms)},
  ymin=0, ymax=45,
  ytick={0,10,20,30,40},
]
%% Full model kappa_n (true)
\addplot[thick,thesisgray,densely dashed,mark=none,domain=0:200,samples=80]
  { (x<5) ? 40 :
    (x<60) ? max(8, 40*exp(-x/18)) :
    max(8, 9 + 0.04*(x-60)) };
\addlegendentry{$\kappa(\mathbf{H})$ (full)};

%% Reduced model kappa (true)
\addplot[thick,thesisblue,mark=none,domain=0:200,samples=80]
  { (x<5) ? 38 :
    (x<60) ? max(6, 38*exp(-x/19)) :
    max(6, 7 + 0.03*(x-60)) };
\addlegendentry{$\kappa(\tilde{\mathbf{H}})$ (reduced)};

%% Upper bound: C_rho * kappa(H)
\addplot[thick,thesisorange,densely dotted,mark=none,domain=0:200,samples=80]
  { min(44,
    ((x<5) ? 40 :
     (x<60) ? max(8, 40*exp(-x/18)) :
     max(8, 9 + 0.04*(x-60)))
    * (1 + ((x<10)?0.02:(x<80)?0.02+(x-10)*0.00085:min(0.19,0.08+(x-80)*0.00092)))
    / (1 - ((x<10)?0.02:(x<80)?0.02+(x-10)*0.00085:min(0.19,0.08+(x-80)*0.00092))) ) };
\addlegendentry{Theorem~1 bound $C_\rho\,\kappa(\mathbf{H})$};

%% kappa = 30 line
\addplot[thesisred,thick,dotted,domain=0:200,samples=2] {30};
\node[thesisred,font=\tiny,anchor=west] at (axis cs:202,31) {30};

\end{groupplot}
\end{tikzpicture}
